\section{Causal RIII search: review and unambiguous certificates}
\label{sec:causal-r3}

The next intake placed three archives in arrival order as reports 26--28:
determinant continuations, causal RIII search, and classical closure resets.
The second archive calls itself ``Report 25'' internally; our collection
number is 27 to avoid colliding with the earlier Garside archive.
Original checksum manifests were verified in scratch before placement,
then omitted under the repository intake rules.  Delivered source and
results remain unchanged apart from the mandated line-ending normalization
of report 28's benchmark CSV.  The retained intake audit records all
original hashes, normalized hashes and source pins.  Local unit reruns pass
35, 7 and 29 tests respectively; these counts do not establish mathematical
soundness or integration with the current recognizer.

\subsection{What the causal argument can give}
Report 27 considers bounded-local reversible moves on a fixed set of sites.
Each move has a footprint containing everything it reads and changes.
Actions on disjoint footprints commute in both directions and preserve
each other's legality.  Terminal reductions also have local footprints.
A move is a \emph{birth} when its footprint is disjoint from the union of
all preceding footprints.  That union is accumulated support, not merely
the footprint of the last move.

Every birth can commute to the front of the trace: it is disjoint from
every earlier action.  The births are pairwise disjoint and can be ordered
canonically.  Remaining actions retain an intersection with the preceding
support.  A reduction may become available sooner after this rearrangement,
in which case the trace is truncated.  Thus one can search a normal form
with a canonical sequence of births followed by connected extensions.
Connectedness of the full dependency graph does not imply that every
chronological prefix was connected.

For RIII on a fixed input diagram, let $D(T)$ be the twelve darts at the
three crossings of a triangular face $T$, and let $\alpha$ be the edge
pairing.  The conservative footprint
\[
 F(T)=D(T)\cup\alpha(D(T)),\qquad |F(T)|\le24
\]
contains every modified pairing.  Inverting the face leaves this footprint
unchanged.  Globally there are $O(N)$ triangular faces in a connected
$N$-crossing spherical diagram.  The number whose footprint meets a set
of $O(k)$ darts is $O(k)$, since only incident crossings and their partners
need inspection.  The report consequently obtains a conditional search
bound
\[
 N^{b+O(1)}(Ck)^k
\]
for traces of length at most $k$ and birth number at most $b$, with an
absolute constant $C$.  A depth-first implementation without a global
memo table uses polynomial workspace in $N,k$.  The delivered abstract
prototype does retain a memo table, so its memory usage should not be
identified with the latter claim.

For example, fixed $b$ and $k=O(\log N/\log\log N)$ give polynomial time;
$b,k=O(\log N)$ give quasi-polynomial time.  These are parameterized
statements.  No argument bounds the parameters uniformly for arbitrary
unknots.  Repeated preprocessing also requires the bound at each episode
selected by the actual deterministic reduction policy; existence of a
different successful sequence of choices is insufficient.

We reran the report's independent breadth-first comparisons in scratch.
The generated audit covers 5,000 six-site systems, with 1,660 witnessed
instances, 3,209 already reducible and 131 without a depth-six witness.
The exhaustive three-site catalog covers 608,400 instances through depth
four, with 171,968 witnessed, 163,800 already reducible and 272,632 without
a witness.  Both reproduce zero discrepancies in the recorded normal-form
checks.  These are abstract Boolean systems, not knot-diagram benchmarks;
they exercise the prototype, not the proposed production adapter.

\subsection{Why immediate inverse pruning needs care}
The proposed RIII adapter suppresses the inverse of the preceding move.
The abstract prototype deliberately does not.  For unrestricted shortest
paths an adjacent inverse pair is removable, but this need not preserve
the birth bound.  Here is an explicit four-site example under precisely
the abstract locality assumptions.  Start with all bits zero:
\begin{center}
\begin{tabular}{llll}
\toprule
Move & Footprint & Toggled bits & Guard\\
\midrule
$A$ & $\{0,1,2,3\}$ & $\{0,1\}$ & $x_2=x_3=0$\\
$B$ & $\{0,2\}$ & $\{2\}$ & $x_0=0$\\
$C$ & $\{1,3\}$ & $\{3\}$ & $x_1=0$\\
\bottomrule
\end{tabular}
\end{center}
Each guard avoids the toggled bits, so every legal move is an involution.
The terminal predicate is $x_2=x_3=1$.  The trace $BC$ succeeds but has
two births.  The trace $AABC$ succeeds with one birth: after $AA$ the
physical state is again zero, while accumulated support contains all four
sites.  After the first $A$, both $B$ and $C$ are disabled, so the inverse
$A$ is essential for this one-birth trace.

Exhaustion of all words through length four gives exactly $AABC$ and
$AACB$ as first-unlocking traces with at most one birth.  Suppressing
adjacent inverses loses both.  The script
\path{data/causal_inverse_audit.py} compares this independent enumeration
with the delivered prototype and retains the source hashes and results.
This is an abstract counterexample to the pruning justification; it is
not an actual RIII counterexample.  It shows that locality alone cannot
justify transplanting that extra pruning rule into the claimed complete
parameterized search.  Physical-state memoization likewise cannot ignore
support and search phase.  The existing bounded local simplifier remains
a heuristic; its failure is not a knottedness certificate.

\subsection{Integrated prerequisite: identify the actual triangular face}
Review also exposed a separate defect in the existing production trace
format.  Recording only the three crossing indices does not uniquely
identify an RIII move.  The two-strand braid $(1,1,-1)$ has the PD code
\[
 ((5,4,0,1),(1,0,2,3),(2,4,5,3)).
\]
Both $(11,4,0)$ and $(9,2,6)$ are legal triangular faces on crossings
$\{0,1,2\}$.  Their inversions give distinct labeled PD codes:
\[
\begin{split}
 ((0,1,2,0),(3,4,1,3),(2,4,5,5)),\\
 ((0,1,1,2),(3,4,4,0),(5,5,3,2)).
\end{split}
\]
The old replay routine selected a face according to traversal order and
therefore could reconstruct the wrong intermediate diagram.  Both moves
preserve knot type; the defect concerns reproducibility of the recorded
sequence, including the legality of its subsequent moves.

New RIII records include a \code{triangle} field containing the face's
three input dart indices.  These indices remain fixed after earlier
crossing deletions.  Replay checks distinct, present crossing indices,
literal integer types, range, face membership and RIII legality.  It
recovers cyclic order from the current pairing, so reordering the three
supplied dart indices cannot accidentally request a different reconnection.
Legacy crossing-only RIII records are accepted when they identify a unique
face, and rejected with an ambiguity error otherwise.  R1/R2 serialization
is unchanged.  This is a certificate correctness fix, not a new recognition
criterion or an asymptotic improvement.

Five new regression tests cover the explicit two-face example with all
six dart orders, invalid certificates, every candidate RIII face in
deterministically generated knot diagrams, legacy unambiguous replay, and
JSON round trips after crossing deletions.  The full integrated suite passes
539 tests in 109.018 seconds in this local run; the log is
\path{data/reidemeister-trace-integrated-tests.txt}.  The clustered search itself
is not yet enabled in production.  Its integration still needs complete
support-aware enumeration, exception-safe rollback, cooperative resource
checks and actual-diagram performance comparisons.  A local search budget
may stop this preprocessing inconclusively; it cannot replace the exact
fallback or establish a general sub-exponential bound.
